\renewcommand{\bibname}{References and Further Reading}
\AddToHookNext{cmd/@makeschapterhead/before}{%
  \phantomsection
  \addcontentsline{toc}{chapter}{References and Further Reading}}
\begin{thebibliography}{99}

\bibitem{bartle-elements}
Robert G. Bartle. \emph{The Elements of Real Analysis}.
Second edition, John Wiley \& Sons, 1976.
\newblock A broad classical companion for Cartesian-space topology,
sequences and function convergence, asymptotic and order notation,
Riemann integration and convergence, finite-dimensional differentiation,
and multivariable integration and change of variables. The manifold and
differential-form development instead uses Munkres, Spivak, Tu, and Lee
as its geometric references.

\bibitem{fis-linear}
Stephen H. Friedberg, Arnold J. Insel, and Lawrence E. Spence.
\emph{Linear Algebra}. Fifth edition, Pearson.
\newblock Chapters 1--2 provide the organizational benchmark for vector
spaces, replacement, bases, linear maps, and their matrix representations.
The analysis main text selects only the finite-dimensional real foundations.

\bibitem{apostol-calculus-two}
Tom M. Apostol. \emph{Calculus, Volume 2}. Second edition, Wiley.
\newblock Further reading for the integration of linear algebra with
multivariable calculus and a structural approach to determinants.

\bibitem{chen-algebra}
Tommy Chen. \emph{Algebra lecture notes}. Author's local revision draft,
consulted September 12, 2026.
\newblock Theorem 5.12.3, ``Existence and uniqueness of the determinant
form,'' supplies the slot-by-slot expansion pattern adapted in Chapter~10.
Section 1.5 guides the concrete permutation discussion in Appendix~\ref{app:determinants}.
Chapter 8 informs the treatment of universal properties in Appendix~\ref{app:tensors},
specialized here to vector spaces over fields.

\bibitem{lee-smooth}
John M. Lee.
\newblock \emph{Introduction to Smooth Manifolds}.
\newblock Graduate Texts in Mathematics 218, second edition, Springer, 2013.
\newblock Primary supplementary technical benchmark for Chapter~13: boundary
charts, tangent spaces, orientation, partitions of unity, integration, and
Stokes. Appendices~\ref{app:submanifolds}--\ref{app:cohomology} compare Chapters 4--8, 17--18, and 20: rank
and slice theorems, embeddings, Lie groups and brackets, cohomology and
integration, and one-parameter subgroups. The strong Whitney theorem
(Theorem 6.19) and de Rham theorem (Theorem 18.14) are stated without proof.
The earlier imported refinement is localized to Theorem~1.15 and the
coordinate-neighborhood construction in the proof of Theorem~2.23, with
the boundary adaptation in Exercise~2.24.

\bibitem{tao-analysis}
Terence Tao. \emph{Analysis I}. Springer.
\newblock Foundations and rational Cauchy constructions; useful for checking
the order of logical dependencies in Chapters~2--4.

\bibitem{abbott-analysis}
Stephen Abbott. \emph{Understanding Analysis}. Second edition, Springer, 2015.
\newblock A companion for motivation, counterexamples, and learning to write
the elementary proofs in Part~I.

\bibitem{ross-analysis}
Kenneth A. Ross. \emph{Elementary Analysis: The Theory of Calculus}.
Second edition, Springer, 2013.
\newblock Elementary rigor and graduated practice with sequences, continuity,
and one-variable calculus.

\bibitem{pugh-analysis}
Charles C. Pugh. \emph{Real Mathematical Analysis}. Second edition,
Springer, 2015.
\newblock Geometric viewpoints on analysis, finite-dimensional topology,
multivariable differentiation, and the geometric volume viewpoint.

\bibitem{knapp-basic-real}
Anthony W. Knapp. \emph{Basic Real Analysis}. Cornerstones, first edition,
Birkh\"auser Boston, 2005.
\newblock A systematic companion for metric and normed spaces,
finite-dimensional differentiation, inverse and implicit functions,
finite partitions of unity, and reduction of substitutions to elementary
coordinate transformations. The author's digital second edition (2016)
is also available from \url{https://www.math.stonybrook.edu/~aknapp/}.

\bibitem{zorich-analysis}
Vladimir A. Zorich. \emph{Mathematical Analysis I--II}. Springer.
\newblock A broader analysis perspective, with geometric interpretations of
multivariable derivatives, integration, and classical vector analysis.

\bibitem{hubbard-vector}
John H. Hubbard and Barbara Burke Hubbard.
\emph{Vector Calculus, Linear Algebra, and Differential Forms:
A Unified Approach}. Fifth edition, Matrix Editions, 2015.
\newblock Geometric companion for interpreting forms as oriented
measurements and connecting them to vector calculus. The publisher's
Chapter~6 sample supplies a selective pedagogical benchmark;
\url{https://matrixeditions.com/UnifiedApproach5thedSamples.html}.

\bibitem{munkres-analysis}
James R. Munkres. \emph{Analysis on Manifolds}. Addison--Wesley, 1991.
\newblock A systematic companion for multivariable differentiation,
primitive local diffeomorphisms, integration, and the route from forms to Stokes.

\bibitem{spivak-manifolds}
Michael Spivak. \emph{Calculus on Manifolds}. W. A. Benjamin, 1965.
\newblock A compact account of the unity of the classical integral theorems;
best read after the more detailed developments here.

\bibitem{tu-manifolds}
Loring W. Tu. \emph{An Introduction to Manifolds}. Second edition,
Springer, 2011.
\newblock Primary pedagogical and organizational benchmark for Chapter~13's
expanded manifold introduction: gradual chart and atlas definitions,
tangent and cotangent spaces, differential forms, orientation, and integration.
Sections 3.10 and 4.1--4.2 also benchmark the distinction between
pointwise covectors, their coordinate basis, and varying form fields.
Appendices~\ref{app:submanifolds}--\ref{app:cohomology} also compare Sections 9 and 11 (submanifolds and rank),
14--16 (vector fields, matrix groups, and Lie algebras), and 24--29
(cohomology, examples, and homotopy). The elementary circle and torus
period proofs avoid importing the book's Mayer--Vietoris machinery.
The exposition here is reconstructed in the manuscript's own notation
and dependency order.

\bibitem{rudin-pma}
Walter Rudin.
\newblock \emph{Principles of Mathematical Analysis}.
\newblock International Series in Pure and Applied Mathematics, third edition,
McGraw--Hill, 1976.
\newblock Further reading for real analysis.

\bibitem{munkres-topology}
James R. Munkres.
\newblock \emph{Topology}.
\newblock Second edition, Prentice Hall, 2000.
\newblock Further reading for the general-topology background used sparingly
in Part~II.

\bibitem{spivak-calculus}
Michael Spivak.
\newblock \emph{Calculus}.
\newblock Fourth edition, Publish or Perish, 2008.
\newblock Further reading for rigorous single-variable calculus.

\bibitem{gessel-lagrange}
Ira M. Gessel. Lagrange inversion.
\emph{Journal of Combinatorial Theory, Series A} 144 (2016), 212--249.
\newblock Coefficient formulations and formal power-series proofs of the
optional analytic refinement; \url{https://arxiv.org/abs/1609.05988}.

\bibitem{sokal-implicit}
Alan D. Sokal. A ridiculously simple and explicit implicit function theorem.
\emph{S\'eminaire Lotharingien de Combinatoire} 61A (2009), Article B61Ad,
21 pages. \url{https://www.mat.univie.ac.at/~slc/s/s61Asokal2.html}.
\newblock Analytic convergence and formal versions of Lagrange inversion.

\bibitem{schwartz-change}
J. Schwartz. The formula for change in variables in a multiple integral.
\emph{The American Mathematical Monthly} 61, no.~2 (1954), 81--85.
\newblock An alternative direct Riemann/Jordan proof using a one-sided
estimate and the inverse map.
\url{https://doi.org/10.1080/00029890.1954.11988420}.

\bibitem{hatcher-topology}
Allen Hatcher. \emph{Algebraic Topology}. Cambridge University Press, 2002.
\newblock Further reading and a benchmark for the optional fundamental-group,
CW-complex, homology, and cohomology previews in Section~13.17;
not a claim of wholesale adaptation. The author's freely available edition
is at \url{https://pi.math.cornell.edu/~hatcher/AT/ATpage.html}.

\bibitem{xiao-topology-notes}
Luciena Xiao Xiao. \emph{MAST31023 Introduction to Algebraic Topology
Lecture Notes}. Supplied lecture-note PDF.
\newblock Supplementary pedagogical benchmark and further reading for the
optional algebraic-topology preview: spaces and maps, homotopy, fundamental
groups, singular chains, and CW characteristic maps and examples. The
manuscript keeps real coefficients and does not adopt the source's full course.

\bibitem{xuewang-manifold-notes}
Xuewang Ai Shuxue (author name romanized from the supplied Chinese notes).
\emph{Weifen Liuxing Biji} [Notes on Smooth Manifolds]. Chinese-language
lecture notes; Section~1, example~7, printed page~4.
\newblock Selective benchmark for Appendix~\ref{app:quotients}'s projective-space proof
architecture: scaling quotient, affine ratio charts, and transitions.
The present Hausdorff proof uses the common isometric-quotient criterion.
This acknowledgment applies to that construction,
not to Appendix~\ref{app:quotients} as a whole. No date or publication venue is supplied.

\bibitem{dummit-foote}
David S. Dummit and Richard M. Foote,
\emph{Abstract Algebra}, 3rd ed., Wiley, 2004.
Sections 10.4 and 11.4--11.5 provide a selective comparison for tensor
products, determinant expansions, and exterior powers.

\end{thebibliography}

\noindent\textbf{Attribution note.}
The principal results used in the manuscript's later development are proved
internally, apart from the external refinement theorem explicitly cited in
Chapter~13. The optional analytic Lagrange inversion theorem in Chapter~11
and the Lebesgue transformation-of-measure preview in Chapter~12 are also
explicitly stated without proof and are not prerequisites. Chapter~9 also contains clearly marked, unproved previews of the
Lipschitz/a.e. and Lebesgue forms of the Fundamental Theorem of Calculus and
of the Cantor function; those statements are not prerequisites for the later
development and carry local references. The references listed here provide
verification and further reading. Section~13.17 contains additional explicitly
unproved preview results, including dominated convergence, Hilbert-space Riesz
representation, Cauchy's integral formula, fundamental-group computations,
Euler-characteristic invariance, cellular--singular homology comparison,
the homological Euler formula, and the maximum
principle; none is a prerequisite. Appendix~\ref{app:submanifolds} states Whitney's embedding
theorem without proof. Appendix~\ref{app:foundations} states ordinary and
transfinite recursion without proof, proves Hartogs's lemma, proves the
equivalence of the Axiom of Choice, the Well-Ordering Theorem, and Zorn's
Lemma from those foundations, and then applies Zorn's Lemma to the
vector-space basis theorem. Appendix~\ref{app:cohomology}
states the Snake and Five Lemmas, long exact sequence theorem,
Mayer--Vietoris, and de Rham theorem without full proofs.
The bibliography entries are not claims that all of the
manuscript's proofs or exposition were adapted from those sources. The
author's algebra notes~\cite{chen-algebra} are acknowledged
locally for the determinant expansion pattern, their influence on the
permutation exposition, and Appendix~\ref{app:tensors}'s proof architecture adapted from
their Chapter~8. These specific adaptations do not describe the origin
of the entire manuscript.

\paragraph{A reading route.}
For a second pass through foundations, pair Tao with Abbott or Ross.
For Part~II, use Pugh or Zorich for geometry and Munkres for a systematic
calculus-on-manifolds treatment. Read Spivak for the compressed unifying
picture after the details are familiar. Tu offers a gradual continuation
into manifolds; Lee supplies a more extensive technical reference. None of
their larger scopes is a prerequisite for these notes.
For the algebraic foundations, compare Chapters~1--2 of
Friedberg--Insel--Spence~\cite{fis-linear}; for determinants and their
calculus role, compare Apostol~\cite{apostol-calculus-two} and Munkres.
Appendices~\ref{app:foundations} and~\ref{app:categories} collect foundational
and categorical language. Appendices~\ref{app:determinants}--\ref{app:cohomology}
develop algebraic extensions, quotient constructions, submanifolds,
Lie theory, and de Rham cohomology along explicit dependency routes.
